\begingroup

\section{Complete independent review of the actual Taylor-unit formal gauge}


\begingroup
\sloppy
\small

\noindent\textbf{Source role.} Complete written proof, with original source bytes and reversible notation/publication mapping retained.\par
\noindent\textbf{Exact revision.} \nolinkurl{5cc96c73e8fabda87d87a5a58184f059361a38210a9696cf498caba911423316}\par
\noindent\textbf{Complete source.} \path{sources/next_edition/unit_gauge_conversions_v1/sources/f1_unit_gauge_review/INDEPENDENT_REVIEW.md}\par
\noindent\textbf{Source SHA-256.} \nolinkurl{5cc96c73e8fabda87d87a5a58184f059361a38210a9696cf498caba911423316}\par\medskip
\subsubsection{Independent review of the actual-unit formal gauge}\label{endpoint-f1_unit_gauge_review--independent-review-of-the-actual-unit-formal-gauge}

Date: 2026-09-13. Read \nolinkurl{ACTUAL_TAYLOR_UNIT_FORMAL_GAUGE.md} completely, UG1--19 and its scope/source paragraphs. The original tau-base note, \nolinkurl{arithmetic_input.tex}, and original exponential comparison were read in full in the adjoining independent scaling work. This review does not edit the source note, run its fixtures, or claim a Lean or analytic-convergence result.

\subsubsection{Disposition and one domain clarification}\label{endpoint-f1_unit_gauge_review--disposition-and-one-domain-clarification}

Accept UG1--18 and the formal construction. The pole inverse has the correct order, the coefficientwise completions are essential and correctly specified, and the explicit retract proves the claimed quasi-isomorphism. The tensor map retains the complete direct Taylor units and recovers the original eta without an assumption that their tensor product is a function of the sum coordinate alone.

UG19's residue identity is correct on the original reflection-stable packet used in the tau-base residue construction. That domain should be explicit here: UG1 otherwise admits any full-order selected packet, while dagger is not automatically a well-defined endomorphism of its quotient. The two equivalent precise typings are:

\begin{itemize}
\tightlist
\item
  use the already retained reflection-stable packet, so dagger maps \(E_h\) to \(E_h\); or
\item
  for arbitrary h, put \(h^\dagger(s)=(-1)^d\overline{h(1-\overline s)}\) (the monic reflected polynomial) and take the first argument from \(E_{h^\dagger}\), so its dagger lies in \(E_h\).
\end{itemize}

No change to the full-order arithmetic source or formal gauge is needed. This is the exact domain of the existing residue arrow, not a new analytic hypothesis. The issue and the tensor-completion convention below were recorded before writing this review.

\subsubsection{Coefficientwise quotient and its inverse}\label{endpoint-f1_unit_gauge_review--coefficientwise-quotient-and-its-inverse}

Let \(A=\mathbb C[s]\), \(A_\nu=\mathbb C[s,\nu^{-1}]\), and \(M=A_\nu/A\). A rational function has a unique decomposition into a polynomial and a proper rational function; two proper representatives differing by a polynomial are equal. Thus the stated section sigma is well defined, including when different denominator powers represent the same rational function. Applying this section to each \((u,t)\) coefficient proves surjectivity in UG5 without assuming a uniform denominator. The constant-nu case gives \(M=0\).

Since \(\gcd(h,\nu)=1\), \(hP/\nu^n\) being polynomial forces \(\nu^n\mid P\). Conversely \(a_n h+b_n\nu^n=1\) gives an inverse class \([a_n P/\nu^n]\). These prove injectivity and surjectivity of \(M_h\) on M, and uniqueness makes its inverse H independent of all Bezout choices. This inverse acts on the additional pole quotient, not on \(E_h\). Differentiation descends to M, but no commutation with H is needed.

Writing \(E=u\partial_s-t\), the exact factorization is

\(D_M=M_h(I+HE)\).

Thus the inverse must be \((I+HE)^{-1}H\), which is precisely UG8. The powers of \(HE\) raise parameter order; at any fixed coefficient only finitely many terms contribute. Both inverse identities follow by finite geometric sums modulo each parameter power, then coefficientwise completion. This does not use a norm-convergent Neumann series or reorder H and the derivative.

\subsubsection{Explicit retraction and formal finite cohomology}\label{endpoint-f1_unit_gauge_review--explicit-retraction-and-formal-finite-cohomology}

Using \(\pi D=D_M\pi\) and \(J D_M=D_M J=I\), one gets

\(KD=\sigma\pi\), \(r^0=I-\sigma\pi\), \(\pi r^1=0\).

The remaining formulas follow by multiplication: \(r^1D=Dr^0\), \(ri=I\), and \(ir=I-dK-Kd\). The degree-zero and degree-one occurrences of K are correctly typed. Thus no assumption about exactness of a completed cohomology functor is replacing the displayed deformation retraction. The mapping-cone description has the same quotient and the contractible identity-cone kernel.

The completed original polynomial complex also retains its expected finite coefficient basis. At the lowest parameter order, divide by h; the \(u\partial_s-t\) term only affects higher orders. Recursive monic division therefore constructs the unique degree-below-d formal remainder and its formal preimage, coefficient by coefficient. The lowest-order nonzero term of a putative kernel would be killed by the injective polynomial multiplier h, so the differential remains injective. Hence its cohomology is the coefficientwise completed free rank-d module. Unbounded s-degrees or pole orders across formal coefficients are allowed exactly as stated in UG4.

\subsubsection{Gauge, special fibre, and scaling}\label{endpoint-f1_unit_gauge_review--gauge-special-fibre-and-scaling}

Direct Leibniz differentiation gives

\(\nu D\nu^{-1}=u\partial_s+h-t-u\nu'/\nu\).

Both the minus sign and factor u are correct. Multiplication by nu in both degrees is an isomorphism only on the explicitly localized complex, with inverse the same rational \(\nu^{-1}\). At \(u=t=0\), both target differentials are h, but this cochain map still multiplies by nu. Its induced map under the original quotient identification is therefore exactly multiplication by \(\upsilon_h\), not identity.

For the composite to the gauge target, an explicit inverse up to homotopy is \(r M_{\nu^{-1}}\); its homotopy is \(M_\nu K M_{\nu^{-1}}\). This also supplies a direct way to tensor the composite itself.

The operator \(L=u\partial_t-s\) commutes with multiplication by nu and its inverse, because these depend only on s. Therefore its commutation with D transfers to \(D^\nu\) with no additional term. The inclusions and gauge map strictly intertwine L and its formal a-adic exponential. It is correct not to demand that the particular proper-part retraction strictly intertwine L. The coefficient matrix in the transported frame \(\nu s^j\) is consequently the same formal \(C_a\); an entire evaluation of that finite matrix does not claim analytic convergence of the rational formal gauge. The note keeps these domains distinct.

\subsubsection{Precise tensor-completion interpretation and homotopy}\label{endpoint-f1_unit_gauge_review--precise-tensor-completion-interpretation-and-homotopy}

The tensor statement is valid with independent parameter pairs and the completed tensor product defined by the cofinal rectangular truncations \((u_i,t_i)^{n_i}\) in the finitely many factors. Each such quotient is the ordinary tensor product over C of its factor truncations. These rectangular filtrations are cofinal with total parameter-adic truncation, so they give the same coefficientwise formal product. Calling every total-degree truncation a literal tensor product would be inaccurate; the rectangular convention is the precise one under which UG7's tensor paragraph operates.

There is no need to infer inverse-limit exactness abstractly. Let \(P_i=i_i r_i\) and use the source's homotopies \(K_i\). For two factors set

\(K_{\mathrm{tot}}=K_1\otimes I+P_1\otimes K_2\),

with the graded convention \((P_1\otimes K_2)(x\otimes y)=(-1)^{\deg x}P_1x\otimes K_2y\). Since P is a cochain map, direct expansion gives

\(dK_{\mathrm{tot}}+K_{\mathrm{tot}}d=(I-P_1)\otimes I+P_1\otimes(I-P_2)\) \(=I-P_1\otimes P_2\).

For k factors the sum has \(P_1,\ldots,P_{j-1}\) before \(K_j\) and identities after it, with these same tensor signs. The terms telescope to \(I-\bigotimes_i P_i\). All maps preserve parameter order, so this identity holds at each compatible rectangular truncation and hence coefficientwise in the completed product. The gauge-conjugated homotopies give the same statement for the targets with \(D_i^\nu\).

Specialization yields \(M_{\upsilon_h}^{\otimes k}\) on the whole original tensor packet. Composing with \([P]\mapsto P(\sum A_i)1\) is exactly eta. All same-sum occupancies remain in the tensor module and no cyclic polynomial formula for the tensor unit is used.

\subsubsection{Residue scope, support, and limits of this review}\label{endpoint-f1_unit_gauge_review--residue-scope-support-and-limits-of-this-review}

With the dagger domain fixed as above, \(g=h(g/h)\) proves UG19: the complete inverse Taylor unit differs from the inverse germ by a multiple of h, so the residue difference is holomorphic. Its transpose and inverse-transpose transports keep the actual source inverse units; they do not assert a positive pairing or a globally invertible extension on every analytic deformed fibre.

Every operation is on the declared linear coefficient fibre, followed by its support-preserving lift. Inverting nu is not inverting the supported zero e. The formal pole contraction does not identify represented zero with tau or pass from tau-base pushforward to the generic skeleton. No specialization of UG8 at a nonzero complex u, no analytic period assertion, no metric estimate, and no tensor-growth bound is certified by this review. The affirmative result is the explicit formal quasi-isomorphism with its actual arithmetic unit and all the displayed kernels retained.

\subsubsection{Final acceptance of the revised note}\label{endpoint-f1_unit_gauge_review--final-acceptance-of-the-revised-note}

The preceding domain-clarification discussion is the historical review record, not an outstanding objection. On 2026-09-13 I reread the revised Section 7 and Section 8, including the full signed tensor homotopy and the explicit original reflection-stable packet restriction. Both comments are resolved in the text. In the tensor discussion above, ``UG7's tensor paragraph'' means Section 7, not the equation labelled UG7.

The exact accepted source artifact is \nolinkurl{ACTUAL_TAYLOR_UNIT_FORMAL_GAUGE.md}, SHA-256:

\nolinkurl{9ded2486160b76bf65f92d44b73cdd0c59a3ca7e5334f88682e10592d5a45b2f}.

Final disposition: accept UG1--19 and the stated formal construction, with no outstanding mathematical correction. Independent parameter pairs and cofinal rectangular truncations make every tensor identity coefficientwise well defined; the displayed Koszul signs give \(dK_{\mathrm{total}}+K_{\mathrm{total}}d=I-\bigotimes_i P_i\). Section 8 now has exactly the domain on which the original dagger and residue-dual arrow are defined, while the preceding gauge proof retains its greater generality. No assumption of analytic convergence, positivity, or global invertibility of the original arithmetic residue unit is added by this acceptance.

This final seal is a proof review. The separately reported exact fixtures and negative controls were not rerun by this reviewer.

\endgroup


\endgroup
